% problem_id: S047-aif3605-N2
% source_id: S047 (Magnetization in the zig-zag layered Ising model and orthogonal polynomials (D. Chelkak, C. Hongler, R. Mahfouf; Ann. Inst. Fourier, Grenoble 74, 6 (2024) 2275-2330; DOI 10.5802/aif.3605))
% source_file: aif3605.pdf
% statement_locator: printed p. 39
% proof_locator: printed p. [39, 40]
% representation: PDF; transcribed from rendered page images (never from the text layer)
% collected_by: carrier rule (the headline theorem was an assembly; this is the step it relies on)
% review_status: PENDING independent review (single delegated judgement)
% status: complete (statement + author proof verbatim + reason + locators)
%
%% problem_id: S-numdam-aif3605-thm-4.4
%% source_id: numdam/aif3605  (D. Chelkak, C. Hongler, R. Mahfouf, "Magnetization in the zig-zag layered
%%            Ising model and orthogonal polynomials", Ann. Inst. Fourier, Grenoble 74, 6 (2024) 2275-2330,
%%            DOI 10.5802/aif.3605)
%% source_file: aif3605.pdf  (host: /disks/sata1/yupeng/human-proof-corpus/source-cache/registry-sources/numdam/aif3605.pdf)
%% candidate: N2
%% target item: Theorem 4.4, section 4.3 ("Boundary magnetic field and the wetting phase transition")
%% statement_locator: printed page 2312 = PDF page 39 (page header "2312   Dmitry CHELKAK, Clement HONGLER & Remy MAHFOUF")
%% proof_locator:     printed pages 2312-2313 = PDF pages 39-40 (proof box ends with "\square" on printed 2313)
%% representation: PDF; transcribed by eye from rendered page images (pdftoppm -r 150 -png for whole pages, plus
%%                 pdftoppm -r 400/-r 600 -png crops for every symbol-level check). The PDF text layer is broken in
%%                 exactly the places that matter here (it prints "+ (1 - r)^{3/2} gamma_{k+n}" for the "- beta_{k+n}"
%%                 of (4.14), it loses the difference between the several zeta/xi/theta variants, and it drops the
%%                 blackboard-bold and diamond glyphs), so pdftotext was used ONLY to locate pages and line numbers.
%%                 No character below was taken from the text layer.
%% status: TRANSCRIPTION ONLY (difficulty judgement recorded separately in N2.json)
%%
%% Transcription notes:
%%  * The unit transcribed is Theorem 4.4 together with its complete proof (printed 2312-2313). Answers to
%%    "which theorem is collected" are in N2.json; the two Remarks 4.5 and 4.6 that follow the proof on the same
%%    page are reproduced as well, marked as FOLLOWING CONTEXT, because they are the immediate reading of (4.14).
%%  * Printed page 2311 = PDF page 38 is reproduced as LEADING CONTEXT only: (4.13) is the list of objects
%%    (q, r, w, xi) in terms of which Theorem 4.4 is stated, and the theorem says "in the setup described above".
%%    Nothing else of printed 2311 is transcribed.
%%  * The published text contains several internal inconsistencies. Each one is kept EXACTLY as printed and is
%%    flagged by a %% [NOTE: ...] line. Nothing has been silently corrected. The four items are:
%%      (a) the last display of the proof prints "+ beta_{k+n}", whereas (4.14) and the identity immediately
%%          above that display both print "- beta_{k+n}"; the identity above is the one that agrees with (4.14);
%%      (b) the same last display prints "c . zeta_0^{k+n-1}", whereas (4.14) prints "gamma_{k+n}" with
%%          gamma_s := c (q^2/r)^s, i.e. c zeta_0^{k+n}; the two differ by zeta_0^{-1};
%%      (c) the normalizing factor is printed "[prod_{k=1}^{2n} cos theta_k]^{-1}" (the source of (1.5) has 2m)
%%          and is said to equal "(1-r)^{-1/2} . (1+q^2)^k" (an exponent m is what the computation needs);
%%      (d) the identity sentence prints w(e^{it}) where the surrounding line has w(e^{i theta}).
%%    See the %% [NOTE: ...] lines in place.
%%  * No passage of the transcribed pages was illegible; there are therefore no [ILLEGIBLE: ...] markers.
%%  * Commas that the journal prints with a preceding thin space (French style, e.g. "[1 - 4a , 1]") are
%%    transcribed with the space kept where it is visible at 600 dpi.

\documentclass[11pt]{article}
\usepackage{amsmath}
\usepackage{amssymb}
\usepackage{amsthm}
\usepackage{mathrsfs}
\usepackage[margin=1in]{geometry}

\theoremstyle{plain}
\newtheorem{theorem}{Theorem}
\newtheorem{remark}{Remark}

\newcommand{\E}{\mathbb{E}}
\newcommand{\Hb}{\mathbb{H}}
\newcommand{\diamondH}{\Hb_{\diamond}}

\begin{document}

%% =====================================================================================
%% --- page 38 (printed 2311) --- LEADING CONTEXT ONLY
%% =====================================================================================

%% [NOTE: everything in this block belongs to printed page 2311 = PDF page 38, i.e. to the page BEFORE the
%%  statement of Theorem 4.4. It is reproduced because Theorem 4.4 says "in the setup described above" and the
%%  setup is exactly the display (4.13) below. The pages 39-40 transcription starts at the next marker.]

\subsection*{4.3. Boundary magnetic field and the wetting phase transition}

In this section we assume that $\theta_k = \theta < \frac{\pi}{4}$ for all $k \geqslant 2$, i.e., that we
work with a fully homogeneous subcritical model but we allow the first interaction constant to have a
different value. This can be trivially reformulated as inducing an additional magnetic field at the
\emph{first column} whose strength $h = 2J_1$ corresponds to $\theta_1$ via (2.1). The main result is the
following theorem which translates the abstract formula (1.5) into the concrete language of
Toeplitz+Hankel determinants. Let
\begin{equation}
\tag{4.13}
\begin{aligned}
& q := \tan\theta < 1 \,, \qquad r := 1 - \frac{\cos^2\theta_1}{\cos^2\theta} \in (-q^2;1) \,, \\
& w(z) := |1 - q^2z| \,, \qquad \xi(z) := \frac{(rz-q^2)(q^2z-1)}{(z-q^2)(q^2z-r)} \,.
\end{aligned}
\end{equation}
Note that $\xi(z)\xi(z^{-1}) = 1$.

%% =====================================================================================
%% --- page 39 (printed 2312) ---
%% =====================================================================================

\medskip
\noindent\textsc{Theorem 4.4.} --- In the setup described above, the following formula holds:
\begin{equation}
\tag{4.14}
M_m = (1-r)^{-3/2}\det\left[\alpha_{k-n} - \beta_{k+n} + (1-r)^{3/2}\gamma_{k+n}\right]_{k,n=0}^{m-1},
\end{equation}
where
\[
\alpha_s := \frac{1}{2\pi}\int_{-\pi}^{\pi} \mathrm{e}^{-is\theta}\, w(\mathrm{e}^{i\theta})\, \mathrm{d}\theta \,, \qquad
\beta_s := \frac{1}{2\pi}\int_{-\pi}^{\pi} \mathrm{e}^{-is\theta}\, \xi(\mathrm{e}^{i\theta})w(\mathrm{e}^{i\theta})\, \mathrm{d}\theta \,,
\]
and $\gamma_s := c\cdot(q^2/r)^s$, $c = (r^2-q^4)r^{-3/2}(r-q^4)^{-1/2}$, if $r > q^2$ and $\gamma_s := 0$
otherwise.

\medskip
\noindent\emph{Proof.} --- Denote $a := \sin^2\theta\cos^2\theta = (q+q^{-1})^{-2}$. The entries of the
Jacobi matrix $J$ (see (1.4)) are given by
\[
b_1 = (1-r)q^{-2}a \,, \quad a_1 = (1-r)^{1/2}a \,; \qquad b_k = 1-2a \,, \quad a_k = a \,, \quad k \geqslant 2 \,.
\]
Let $\varrho_k := (1-r\delta_{k,0})^{1/2}$, where $\delta_{k,0}$ is the Kronecker delta. The continuous
spectrum of $J$ has multiplicity $1$ and equals to $[1-4a\,,1]$. The generalized eigenfunctions are
\[
\psi_k(\zeta) := \varrho_k^{-1}\cdot\left[\zeta^k - \xi(\zeta)\zeta^{-k}\right] \,, \qquad
\lambda(\zeta) := 1 - a\cdot(2+\zeta+\zeta^{-1}) \,,
\]
where $\zeta = \mathrm{e}^{i\theta}$ and $\theta \in [0,\pi]$. The coefficient $\xi(\zeta)$ should satisfy the
condition $(b_1-\lambda(\zeta))\psi_0(\zeta) = a_1\psi_1(\zeta)$ which leads to the formula (4.13). The matrix
$J$ also has the eigenvalue
\[
\lambda(\zeta_0) = \frac{(1-r)(r-q^4)}{r(1+q^2)^2} \in (0,1-4a) \qquad \text{if } \zeta_0 := q^2/r < 1
\]
since $\xi(\zeta_0) = 0$. Note that
\begin{align*}
\frac{\varrho_k\varrho_n}{2\pi}\int_0^{\pi} \psi_n(\mathrm{e}^{-i\theta})\psi_k(\mathrm{e}^{i\theta})\,\mathrm{d}\theta
& = \frac{1}{2\pi \mathrm{i}}\oint_{|\zeta|=1} \left[\zeta^{k-n} - \xi(\zeta^{-1})\zeta^{k+n}\right]\frac{\mathrm{d}\zeta}{\zeta} \\
& = (1-r\delta_{k+n,0})\cdot\delta_{k,n} \;-\; c_0\zeta_0^{k+n-1} \,,
\end{align*}
where $c_0 = 0$ if $r \leqslant q^2$ and
\[
c_0 := \operatorname{res}_{z=\zeta_0}\xi(z^{-1}) = \frac{q^2(1-r)(r^2-q^4)}{r^2(r-q^4)} \qquad \text{if } r > q^2 \,.
\]
Thus, the spectral decomposition of the basis vector $e_n = (\delta_{k,n})_{k\geqslant 0}$ reads as
\[
\delta_{k,n} = \frac{1}{2\pi}\int_0^{\pi} \psi_n(\mathrm{e}^{-i\theta})\psi_k(\mathrm{e}^{i\theta})\,\mathrm{d}\theta
+ \varrho_n^{-1}c_0\zeta_0^{n-1}\cdot\psi_k(\zeta_0) \,.
\]

%% =====================================================================================
%% --- page 40 (printed 2313) ---
%% =====================================================================================

%% [NOTE: the sentence below is printed with "\mathrm{e}^{it}" inside an integral over the circle; the
%%  corresponding object everywhere else on pages 2312-2313 is w(e^{i\theta}). Printed as it stands.]

Since $\lambda(\mathrm{e}^{i\theta}) = (1+q^2)^{-2}(w(\mathrm{e}^{it}))^2$, this gives the identity
\begin{align*}
\varrho_k\varrho_n\langle e_k, J^{1/2}e_n\rangle
& = \frac{\varrho_k\varrho_n}{2\pi}\int_0^{\pi} \psi_n(\mathrm{e}^{-i\theta})\psi_k(\mathrm{e}^{i\theta})
    \frac{w(\mathrm{e}^{i\theta})\,\mathrm{d}\theta}{1+q^2} \;+\; c_0\zeta_0^{k+n-1}(\lambda(\zeta_0))^{1/2} \\
& = \varrho_k\varrho_n\cdot\left[(1+q^2)^{-1}(\alpha_{k-n}-\beta_{k+n}) + c_0(\lambda(\zeta_0))^{1/2}\zeta_0^{k+n-1}\right].
\end{align*}

%% [NOTE: the two internal inconsistencies (a) and (b) of the transcription notes both sit in the display that
%%  follows. The source prints "+ \beta_{k+n}" where (4.14) and the identity just above print "- \beta_{k+n}",
%%  and the source prints the exponent "k+n-1" where (4.14) asks for gamma_{k+n} = c (q^2/r)^{k+n}. Both are
%%  reproduced exactly as printed, not corrected.]

It remains to note that the normalizing factor $\left[\prod_{k=1}^{2n}\cos\theta_k\right]^{-1}$ in (1.5) equals
to $(1-r)^{-1/2}\cdot(1+q^2)^k$ and hence (note also the two factors $\varrho_0 = (1-r)^{-1/2}$ in the first row
and the first column of the matrix $J^{1/2}$)

%% [NOTE: in the source the multiplication dot between "c" and "\zeta_0" is printed with an unusually heavy
%%  horizontal stroke; at 600 dpi it reads as an ordinary (if oddly inked) "\cdot". Nothing else is meant.]

\[
M_m = (1-r)^{-3/2}\det\left[\alpha_{k-n} + \beta_{k+n} + (1-r)^{3/2}c\cdot\zeta_0^{k+n-1}\right]_{k,n=0}^{m-1},
\]
where
\[
c := \frac{r(1+q^2)c_0(\lambda(\zeta_0))^{1/2}}{q^2(1-r)^{3/2}} = \frac{r^2-q^4}{r^{3/2}(r-q^4)^{1/2}}
\]
as claimed. \hfill$\square$

%% -------------------------------------------------------------------------------------
%% FOLLOWING CONTEXT (same page, printed 2313; not part of the proof of Theorem 4.4)
%% -------------------------------------------------------------------------------------

\medskip
\noindent\emph{Remark 4.5} (free boundary conditions). --- One can pass to the limit $r \to 1^-$ (which
corresponds to $J_1 \to 0^+$) in the formula (4.14) since $\alpha_s = \alpha_{-s} = \beta_s + O(1-r)$ and
$\alpha_0 = \beta_0 + O((1-r)^2)$ as $r \to 1^-$. (It is also not hard to adapt the proofs of Theorems 1.1 and
4.4 for this setup.) In particular, one can easily see that
\[
\E^{+,0}_{\diamondH}\left[\sigma_{\left(-\frac{5}{2},0\right)}\right] = (1-q^4)^{1/2} \,, \qquad r = 1 \,,
\]
where the sign ``$+$'' in the superscript indicates the boundary conditions at infinity and $0$ stands for the
value of the magnetic field $h$ at the vertical boundary (free boundary conditions). Note that $M_1$ does not
vanish at $h = 0$ provided that $q < 1$: the ``$+$'' boundary conditions at infinity break the spin-flip
symmetry.

\medskip
\noindent\emph{Remark 4.6} (wetting phase transition). --- In fact, one can analytically continue the
right-hand side of (4.14) to negative values of $(1-r)^{1/2}$. According to [26, 53], this corresponds to a
wetting phase transition. Informally speaking, for small negative values $-h$ of the boundary magnetic field,
the interface separating ``$+$'' boundary conditions at infinity from ``$-$'' ones on the imaginary line
$i\mathbb{R}$ touches the boundary infinitely often and the ``$+$'' phase dominates in the bulk of the
half-plane, while for big negative values $-h$ this interface ``breaks away'' from $i\mathbb{R}$ and the
``$-$'' phase dominates in the bulk. For instance, one should have
\[
\E^{+,-h}_{\diamondH}\left[\sigma_{\left(-\frac{5}{2},0\right)}\right]
= -|1-r|^{3/2}(\alpha_0-\beta_0) + \gamma_0
= 2\gamma_0 - \E^{+,h}_{\diamondH}\left[\sigma_{\left(-\frac{5}{2},0\right)}\right]
\]

%% -------------------------------------------------------------------------------------
%% END OF UNIT. Printed page 2314 = PDF page 41 continues with the discussion of Remarks 4.5-4.6
%% ("...provided that h is small enough. Due to Theorem 4.4, the mismatch 2 gamma_0 disappears ..."),
%% which uses (4.14); it is not part of the proof and is not transcribed here.
%% -------------------------------------------------------------------------------------

\end{document}
